% Generated by scripts/gen_blueprint.py from blueprint/nodes/03-hermite-lindemann.toml. Do not edit.

\chapter{Hermite–Lindemann and Lindemann–Weierstrass}
\label{chap:03-hermite-lindemann}

The Lean proof of Lindemann--Weierstrass is ported from Mathlib pull request \#28013 by Yuyang Zhao, which is not merged; it depends on Mathlib alone. Hermite--Lindemann is derived from it, in two equivalent forms. The four exponentials development uses Hermite--Lindemann to make one generator of a field transcendental.

% Prove2Me: Schanuel.lindemann_weierstrass -- https://prove2.me/theorems/afb75b5a-3ca6-4700-9c55-ebd4e3ecc762
% Statement written from the Lean source: part (b), algebraicIndependent_exp, is not on the page, which states part (a)
\begin{theorem}[Lindemann–Weierstrass]
\label{thm:lindemann_weierstrass}
\lean{linearIndependent_exp, algebraicIndependent_exp}
\leanok
\begin{enumerate}
\item[(a)] Let $\alpha_1,\dots,\alpha_n$ be pairwise distinct algebraic numbers. Then $e^{\alpha_1},\dots,e^{\alpha_n}$ are linearly independent over $\Qbar$: if $\beta_1,\dots,\beta_n\in\Qbar$ are not all zero, then $\sum_i\beta_ie^{\alpha_i}\neq0$.
\item[(b)] If algebraic numbers $\alpha_1,\dots,\alpha_n$ are linearly independent over $\Q$, then $e^{\alpha_1},\dots,e^{\alpha_n}$ are algebraically independent over $\Qbar$.
\end{enumerate}
(The Lean statements take families indexed by any type; in (b) they ask for linear independence over $\N$, which for a family in a $\Q$-vector space is the same.)

{\small\emph{Source:} Lindemann \cite{Lin1882}, Weierstrass \cite{Wei1885}; see \cite[Theorem 1.4]{Bak75}.}
\end{theorem}

\begin{proof}
\leanok
(a) is the proof of the Mathlib pull request. (b) applies (a) to the numbers $\sum_ie_i\alpha_i$, $e\in\N^{n}$, which are pairwise distinct.
\end{proof}

% Prove2Me: DiazModulus.hermite_lindemann_holds -- https://prove2.me/theorems/fdc68131-2e60-4005-9489-8758a2174325
\begin{theorem}[Hermite–Lindemann]
\label{thm:hermite_lindemann_holds}
\lean{Diaz.hermite_lindemann_holds}
\leanok
If $a$ is a non-zero algebraic number, then $e^{a}$ is transcendental.

{\small\emph{Source:} Hermite \cite{Her1873} for $a=1$, Lindemann \cite{Lin1882} in general; see \cite[Theorem 1.4]{Bak75}.}
\end{theorem}

\begin{proof}
\leanok
\uses{thm:lindemann_weierstrass}
Apply Theorem~\ref{thm:lindemann_weierstrass}(a) to $\alpha=(a,0)$ and $\beta=(1,-e^{a})$.
\end{proof}

% Prove2Me: Schanuel.hermite_lindemann -- https://prove2.me/theorems/f88819a5-647b-4cbb-854b-5c8bbfe23d70
% Statement written from the Lean source: the page states the equivalent direct form, e^a transcendental for algebraic a ≠ 0
\begin{theorem}[Hermite–Lindemann, logarithmic form]
\label{thm:hermite_lindemann}
\lean{Diaz.hermite_lindemann}
\leanok
Let $u\in\C$, $u\neq0$, with $e^{u}$ algebraic. Then $u$ is transcendental.

{\small\emph{Source:} Hermite \cite{Her1873}, Lindemann \cite{Lin1882}; see \cite[Theorem 1.4]{Bak75}.}
\end{theorem}

\begin{proof}
\leanok
\uses{thm:lindemann_weierstrass}
Apply Theorem~\ref{thm:lindemann_weierstrass}(a) to $\alpha=(u,0)$ and $\beta=(1,-e^{u})$. This is the form the rest of the library uses; Theorem~\ref{thm:hermite_lindemann_holds} is its contrapositive. (The accepted Prove2Me proof derives it from Theorem~\ref{thm:hermite_lindemann_holds}; the library's proof, shown here, does not.)
\end{proof}
